% !TeX root = ../navier-stokes-manuscript.tex
% ============================================================
% 1.8 Dead, Jump, Blown, and Restart
% ============================================================
\subsection{Dead, Jump, Blown, and Restart}\label{sub:DeadJumpBlownAndRestart}

The derivative tower turns a vague loss of smoothness into a failure we can picture inside the fluid.
At every regular point, each response belongs to one surrounding velocity field and satisfies the equation opened at that rung.
Near an alleged endpoint, one portion may go still while pressure and viscous curvature continue to demand change from it.
Two neighbouring approaches may carry incompatible finite values into the same location.
A shrinking pulse or vortex may drive one response beyond every finite size in the time remaining.
These are the three physical scenes developed below.

\divider{Dead}

Imagine one point of the fluid stopping dead in its tracks.
Its velocity may be zero.
Its acceleration may be zero too, and so may its jerk or any other particular rung.
None of those zeroes proves a failure by itself.
The question is whether the fluid around that point is still demanding that something change.

At the base of the tower, that demand is visible in
\[
  D_tu=-\nabla p+\nu\Delta u.
\]
The left side is the response carried by the parcel.
The right side is what pressure and viscous curvature in the surrounding field require from it.
For a proposed extension \(Q=(\widetilde u,\widetilde p)\), set \(\widetilde D_t:=\partial_t+\widetilde u\cdot\nabla\).
If the proposal assigns
\[
  \widetilde D_t\widetilde u=0
\]
while
\[
  -\nabla\widetilde p+\nu\Delta\widetilde u\neq0
\]
at that point, then the point has stopped while its fluid has not stopped asking it to move.
The same picture applies to a material packet \(A(t)\) carried by \(\widetilde u\).
Its momentum response can go dead while the surrounding traction remains nonzero:
\[
  \frac{d}{dt}
  \int_{A(t)}\widetilde u(x,t)\,dx
  =0,
  \qquad
  \int_{\partial A(t)}
  \left[
    -\widetilde pI
    +
    \nu\bigl(\nabla\widetilde u+\nabla\widetilde u^{\mathsf T}\bigr)
  \right]n\,dS
  \neq0.
\]
The packet then refuses the net pressure and viscous response transmitted through its boundary.

Now carry the test up the tower.
Write
\[
  \widetilde V_n:=\partial_t^n\widetilde u,
  \qquad
  \widetilde Q_n:=\partial_t^n\widetilde p.
\]
The row from \cref{sub:SingularityAndTheDerivativeTower} says that the next response must be
\[
  \widetilde V_{n+1}
  =
  \nu\Delta\widetilde V_n
  -
  \sum_{k=0}^{n}
  \binom{n}{k}
  (\widetilde V_k\cdot\nabla)\widetilde V_{n-k}
  -
  \nabla\widetilde Q_n.
\]
Velocity can go to zero, and every chosen rung can go to zero.
But when the surrounding pressure, spatial curvature, and lower responses make the right side nonzero, the rung immediately above them has to remain alive.
If the proposal instead lets the whole local tower go dead, then no response remains to carry information from the surrounding fluid into the point or packet.

Zero can also be completely lawful.
The zero solution answers its surrounding field exactly, a smooth steady solution can have every positive Eulerian time rung equal to zero, and a changing response can pass through zero at an instant.
Those cases remain alive because the equation asks for zero too.
Dead begins only when the responses go still while some row of the full equation continues to demand change.

Define the mismatch in the proposed \(n\)-th row by
\[
  \mathfrak D_n[Q]
  :=
  \widetilde V_{n+1}
  +
  \sum_{k=0}^{n}
  \binom{n}{k}
  (\widetilde V_k\cdot\nabla)\widetilde V_{n-k}
  +
  \nabla\widetilde Q_n
  -
  \nu\Delta\widetilde V_n.
\]
Every classical Navier--Stokes region has \(\mathfrak D_n[Q]=0\).
If \(Q\) is classical and \(\mathfrak D_n[Q](z_0)\neq0\), the equation already fails at \(z_0\), and continuity carries that mismatch into a spacetime neighbourhood of \(z_0\).
If \(Q\) is not pointwise classical but has enough local regularity that every term in \(\mathfrak D_n[Q]\) defines a distribution on a neighbourhood \(\mathcal O\), the witness is
\[
  \mathfrak D_n[Q]
  \neq
  0
  \quad
  \text{in }\mathcal D'(\mathcal O;\mathbb R^3).
\]
Some vector-valued test function then detects it:
\[
  \left\langle
    \mathfrak D_n[Q],\phi
  \right\rangle
  \neq0
\]
for some \(\phi\in C_c^\infty(\mathcal O;\mathbb R^3)\).
That is the mathematical witness of the physical Dead scene: the point or packet has stopped participating, one of the fluid's required properties has failed there, and the proposed field has left the Navier--Stokes class.

\divider{Jump}

Now imagine the response stays finite, yet neighbouring approaches to the same spacetime point refuse to produce one value.
For a temporal, spatial, or mixed rung
\[
  J_{n,\alpha}
  :=
  \partial_t^n\partial_x^\alpha u,
\]
suppose two sequences of regular points \(z_m^+\to z_*\) and \(z_m^-\to z_*\) satisfy
\[
  \left|
    J_{n,\alpha}(z_m^+)
    -
    J_{n,\alpha}(z_m^-)
  \right|
  \geq
  \varepsilon_0>0
\]
along a subsequence.
No single continuous value can join the two approaches at \(z_*\).
The field has Jump at that rung.

The planar Euler sheet is the lowest visible version.
The same one-fluid packet approaches its internal surface from below with velocity \(U_-e_1\) and from above with velocity \(U_+e_1\).
Its missing intermediate velocities prevent the two neighbouring sides from becoming one continuous Navier--Stokes response, and fixed viscosity detects that loss through the distributional curvature of the jump.

The same slip has an overlapping Blown reading when it is approached through thinner continuous layers.
If the velocity difference \(\Delta U=U_+-U_-\) is spread across thickness \(\varepsilon\), its shear has the scale
\[
  \frac{|\Delta U|}{\varepsilon}.
\]
The two velocities stay finite while this gradient diverges as \(\varepsilon\downarrow0\), and the zero-thickness sheet carries the limiting shear as a surface distribution.

A higher Jump can leave velocity continuous.
The two approaches may agree on \(u\) while disagreeing on \(\nabla u\), on a fixed-coordinate time rung \(V_n\), or on the combination of mixed derivatives that forms material acceleration or jerk.
Oscillation can create the same failure of the Cauchy test without producing two simple one-sided traces.
The decisive fact is that one part of the field can no longer be joined continuously to the values carried toward it by neighbouring parts.

Jump can overlap with Dead.
Assigning one of the competing values at the endpoint may create a finite equation mismatch, while leaving the value unassigned still blocks smooth continuation.
The first statement concerns participation in the equation, and the second concerns arrival at one continuous field.

\divider{Blown}

The third scene begins when concentration drives the response upward inside a shrinking region.
Imagine a smooth pulse or vortex structure being worked into a smaller region while its velocity difference is taken across a still shorter distance.
The velocity can remain finite while a gradient, acceleration, jerk, or higher mixed response grows beyond every finite height near the concentrating region.
At a point this appears as
\[
  \limsup_{(x,t)\to(x_*,T)}
  |J_{n,\alpha}(x,t)|
  =\infty,
\]
and on a neighbourhood it can appear as
\[
  \limsup_{t\uparrow T}
  \lVert J_{n,\alpha}(\cdot,t)\rVert_{L^p(B_r(x_*))}
  =\infty.
\]
This is Blown at the indicated rung.

A simple scale picture shows why finite velocity does not protect the tower.
Let velocity increments and length scales satisfy
\[
  a_k=2^{-k},
  \qquad
  \ell_k=4^{-k}.
\]
Then
\[
  \sum_{k=1}^\infty a_k<\infty,
  \qquad
  \frac{a_k}{\ell_k}=2^k\longrightarrow\infty.
\]
The total velocity change stays finite while the corresponding gradient scale diverges.
This is a kinematic picture of how Blown can first appear above velocity.
No Navier--Stokes solution realizing this sequence has been constructed here.

The critical scaling gives the concentration a second face.
A structure of length \(\ell\) and velocity height comparable to \(\ell^{-1}\) keeps its \(L^3\) scale unchanged while its \(L^2\) energy contribution is comparable to \(\ell\).
Its height is critical because shrinking the structure does not shrink that \(L^3\) reading.
At the same time, its parabolic duration is comparable to \(\ell^2\).
For scales \(\ell_k=2^{-k}\), the durations
\[
  \ell_k^2=4^{-k}
\]
fit inside a finite total time.
A pulse or vortex can thus be pictured passing through successively finer critical-scale structures while using less energy and less time at each stage.
This Zeno picture is allowed by the scaling gap in the known estimates; it is not a constructed singular solution.

Blown can meet the other two scenes.
A divergent oscillation can lose finite size and a unique approach together, while a proposed truncation of that divergence can create an equation mismatch.
The overlap matters because these are candidate general pictures of a failure inside the fluid.
No exhaustiveness or exclusivity is claimed for the three scenes.

In all three general pictures, some portion of the fluid no longer responds to information from its neighbours.
Neighbour information includes the global pressure response, because incompressibility generates pressure from the whole velocity field and a distant change can alter the response demanded at the portion under examination.
Dead loses the demanded equation response, Jump loses a common value across neighbouring approaches, and Blown loses finite control as the response concentrates.
Each picture must still be connected to the actual Navier--Stokes solution before it becomes a proof of class exit or failed continuation.

\divider{Restart}

The same solution also contains its own fresh initial slices whenever it is regular.
Choose a regular time \(t_0<T_*\) and set
\[
  v(x,s):=u(x,t_0+s),
  \qquad
  q(x,s):=p(x,t_0+s).
\]
Direct substitution gives
\[
  \partial_sv+(v\cdot\nabla)v+\nabla q
  =
  \nu\Delta v,
  \qquad
  \nabla\cdot v=0,
  \qquad
  v(\cdot,0)=u(\cdot,t_0).
\]
The later regular slice \(u(\cdot,t_0)\) has become new initial data for the same evolution.
Local uniqueness joins the restarted solution to the original one wherever their time intervals overlap.
Changing the time origin changes no velocity, pressure, viscosity, or material parcel in the fluid.
This is the restart that a continuation-grade bound makes long enough to cross a candidate endpoint.
